% TOP_PROBLEM: {"problemId": "problem.aspherical-manifolds-admit-no-metric-of-positive-scalar-curvature", "canonicalTitle": "Aspherical Manifolds Admit No Metric of Positive Scalar Curvature", "releaseRank": 123, "primaryDomain": "geometry_topology", "primaryDomainLabel": "Geometry & topology", "displayStatus": "open_with_solved_subcases", "releaseStatus": "open_with_solved_subcases"}
% !TeX program = lualatex
% Definition and sources only; this file is not an LLM solution attempt.
\documentclass[11pt]{article}
\usepackage[margin=25mm]{geometry}
\usepackage{fontspec}
\setmainfont{Latin Modern Roman}
\usepackage{amsmath,amssymb,mathtools}
\usepackage{unicode-math}
\setmathfont{Latin Modern Math}
\usepackage{xurl,hyperref}
\hypersetup{colorlinks=true,urlcolor=blue}
\setcounter{secnumdepth}{0}
\setlength{\parindent}{0pt}
\setlength{\parskip}{5pt}
\setlength{\emergencystretch}{3em}
\begin{document}
\section*{\texorpdfstring{Aspherical Manifolds Admit No Metric of Positive Scalar Curvature}{Aspherical Manifolds Admit No Metric of Positive Scalar Curvature}}\label{123}
\subsection{Definitions and mathematical statement}
A closed manifold is compact without boundary, and asphericity means its universal cover is contractible. The scalar curvature of a Riemannian metric is the trace of its Ricci curvature. The target is
\[
 M\text{ closed and aspherical}\quad\Longrightarrow\quad
 \nexists g\text{ on }M\text{ with }\operatorname{Scal}_g(x)>0\text{ for every }x.
\]
Strict pointwise positivity is intended. A prohibition only on positive \emph{sectional} curvature is weaker, and additional spin or fundamental-group assumptions would restrict the conjecture.
\par\smallskip\textit{Formulation and concept references:} \href{https://arxiv.org/html/2312.04698}{[S1]}.

\subsection{Short English statement}
Does asphericity rule out every everywhere-positive scalar-curvature metric on a closed manifold?
\subsection{Sources}
\begin{itemize}
\item[S1]Simone Cecchini, Jianchun Chu, and Jintian Zhu, Positive scalar curvature metrics and aspherical summands, arXiv:2312.04698v3 (2025), introduction.
\url{https://arxiv.org/html/2312.04698}.
\textit{Formulation source.}
\end{itemize}
\end{document}
